\documentclass[10pt,a4paper]{article}
\usepackage[spanish,es-nodecimaldot]{babel}
\usepackage[margin=20mm]{geometry}
\usepackage{microtype,booktabs,tabularx,xcolor,fancyhdr}
\usepackage{pgfplots}
\usepackage[colorlinks=true,allcolors=azul]{hyperref}
\pgfplotsset{compat=1.18}

\definecolor{azul}{HTML}{163F6A}
\definecolor{verde}{HTML}{007960}
\definecolor{tinta}{HTML}{172B39}
\definecolor{veladura}{HTML}{EFF5F6}
\color{tinta}
\setlength{\parindent}{0pt}
\setlength{\parskip}{5pt}
\setlength{\headheight}{14pt}
\pagestyle{fancy}
\fancyhf{}
\fancyhead[L]{\small\textcolor{azul}{EDGE-DB · BASES DE DATOS EN LA NUBE}}
\fancyhead[R]{\small Universidad Politécnica de Chiapas}
\fancyfoot[L]{\small Caso de estudio · noviembre de 2026}
\fancyfoot[R]{\small\thepage}
\renewcommand{\headrulewidth}{0.3pt}
\newcommand{\entrada}[1]{\vspace{4pt}{\large\bfseries\color{azul}#1}\par}
\newcommand{\clave}[1]{\texttt{#1}}
% Generado por results/analizar.py a partir del ensayo indicado.
\newcommand{\EdgeDato}[1]{\csname edge-#1\endcsname}
\expandafter\def\csname edge-pilot_runs\endcsname{6}
\expandafter\def\csname edge-final_runs\endcsname{70}
\expandafter\def\csname edge-attempts\endcsname{80}
\expandafter\def\csname edge-failed\endcsname{1}
\expandafter\def\csname edge-superseded_completed\endcsname{3}
\expandafter\def\csname edge-low_p99\endcsname{64}
\expandafter\def\csname edge-total_operations\endcsname{2750}
\expandafter\def\csname edge-idle_samples\endcsname{62}
\expandafter\def\csname edge-idle_duration\endcsname{64.3}
\expandafter\def\csname edge-idle_ram\endcsname{256.1}
\expandafter\def\csname edge-idle_rss\endcsname{73.7}
\expandafter\def\csname edge-run_ram_min\endcsname{188.6}
\expandafter\def\csname edge-run_rss_max\endcsname{368.8}
\expandafter\def\csname edge-run_temp_min\endcsname{48.3}
\expandafter\def\csname edge-run_temp_max\endcsname{68.8}
\expandafter\def\csname edge-run_cpu_max\endcsname{99.8}
\expandafter\def\csname edge-run_swap_max\endcsname{70.2}
\expandafter\def\csname edge-run_read_max\endcsname{7303168}
\expandafter\def\csname edge-run_write_max\endcsname{93696000}
\expandafter\def\csname edge-trace_samples\endcsname{719}
\expandafter\def\csname edge-boot_check\endcsname{coincidió}
\expandafter\def\csname edge-kernel_oom\endcsname{sin coincidencias}
\expandafter\def\csname edge-flags\endcsname{0x0}
\expandafter\def\csname edge-window\endcsname{5}
\expandafter\def\csname edge-warm\endcsname{1}
\expandafter\def\csname edge-cool\endcsname{1}
\expandafter\def\csname edge-read_mib\endcsname{7.0}
\expandafter\def\csname edge-write_mib\endcsname{89.4}
\expandafter\def\csname edge-q1_1000_1_p95\endcsname{85.6}
\expandafter\def\csname edge-q1_1000_1_qps\endcsname{14.0}
\expandafter\def\csname edge-q1_10000_1_p95\endcsname{539.8}
\expandafter\def\csname edge-q1_10000_1_qps\endcsname{1.8}
\expandafter\def\csname edge-q1_10000_2_p95\endcsname{552.6}
\expandafter\def\csname edge-q1_10000_2_qps\endcsname{3.6}
\expandafter\def\csname edge-q1_10000_5_p95\endcsname{797.3}
\expandafter\def\csname edge-q1_10000_5_qps\endcsname{7.6}
\expandafter\def\csname edge-q1_50000_1_p95\endcsname{3115.0}
\expandafter\def\csname edge-q1_50000_1_qps\endcsname{0.3}
\expandafter\def\csname edge-q2_1000_1_p95\endcsname{160.6}
\expandafter\def\csname edge-q2_1000_1_qps\endcsname{20.4}
\expandafter\def\csname edge-q2_10000_1_p95\endcsname{220.6}
\expandafter\def\csname edge-q2_10000_1_qps\endcsname{6.4}
\expandafter\def\csname edge-q2_10000_2_p95\endcsname{246.9}
\expandafter\def\csname edge-q2_10000_2_qps\endcsname{11.6}
\expandafter\def\csname edge-q2_10000_5_p95\endcsname{473.5}
\expandafter\def\csname edge-q2_10000_5_qps\endcsname{21.2}
\expandafter\def\csname edge-q2_50000_1_p95\endcsname{313.3}
\expandafter\def\csname edge-q2_50000_1_qps\endcsname{3.4}
\expandafter\def\csname edge-50000-q1-nodes\endcsname{Aggregate, Incremental Sort, Nested Loop, Gather Merge, Sort, Seq Scan, Function Scan}
\expandafter\def\csname edge-50000-q1-sort\endcsname{external merge}
\expandafter\def\csname edge-50000-q1-temp\endcsname{3532}
\expandafter\def\csname edge-50000-q1-sharedread\endcsname{0}
\expandafter\def\csname edge-50000-q1-sharedhit\endcsname{3938}
\expandafter\def\csname edge-50000-q1-ms\endcsname{3179.7}
\expandafter\def\csname edge-50000-q2-nodes\endcsname{Limit, Gather Merge, Sort, Seq Scan}
\expandafter\def\csname edge-50000-q2-sort\endcsname{top-N heapsort}
\expandafter\def\csname edge-50000-q2-temp\endcsname{0}
\expandafter\def\csname edge-50000-q2-sharedread\endcsname{0}
\expandafter\def\csname edge-50000-q2-sharedhit\endcsname{3950}
\expandafter\def\csname edge-50000-q2-ms\endcsname{305.6}
\expandafter\def\csname edge-50000-q1-jit-off-ms\endcsname{3057.4}
\expandafter\def\csname edge-50000-q1-jit-functions\endcsname{17}
\newcommand{\EdgeRows}{%
1,000 & 1 & Q1 & 62 & 85.6 [82.3, 266.3] & 290.3 & 14.0 \\
1,000 & 1 & Q2 & 84 & 160.6 [57.8, 216.3] & 304.7 & 20.4 \\
5,000 & 1 & Q1 & 14 & 469.3 [252.8, 501.5] & 469.3 & 3.6 \\
5,000 & 1 & Q2 & 10 & 115.3 [112.8, 134.2] & 321.3 & 10.8 \\
10,000 & 1 & Q1 & 9 & 539.8 [505.6, 540.4] & 539.8 & 1.8 \\
10,000 & 1 & Q2 & 25 & 220.6 [199.2, 395.3] & 397.3 & 6.4 \\
10,000 & 2 & Q1 & 18 & 552.6 [541.5, 565.9] & 552.6 & 3.6 \\
10,000 & 2 & Q2 & 44 & 246.9 [226.3, 269.5] & 302.1 & 11.6 \\
10,000 & 5 & Q1 & 37 & 797.3 [776.7, 839.4] & 862.5 & 7.6 \\
10,000 & 5 & Q2 & 92 & 473.5 [335.1, 525.1] & 681.8 & 21.2 \\
25,000 & 1 & Q1 & 2 & 1479.4 [1476.3, 1482.9] & 1479.4 & 0.6 \\
25,000 & 1 & Q2 & 14 & 329.4 [329.0, 410.4] & 329.4 & 3.0 \\
50,000 & 1 & Q1 & 6 & 3115.0 [3033.2, 3450.3] & 3115.0 & 0.3 \\
50,000 & 1 & Q2 & 15 & 313.3 [289.5, 461.2] & 313.3 & 3.4 \\
}
\newcommand{\EdgeDatasetRows}{%
1,000 & 3,023 & 598 & 0.35 & 0.95 \\
5,000 & 14,949 & 3,004 & 1.77 & 3.53 \\
10,000 & 29,939 & 5,960 & 3.54 & 6.64 \\
25,000 & 74,769 & 14,972 & 8.86 & 15.46 \\
50,000 & 149,760 & 30,171 & 17.73 & 30.99 \\
}
\newcommand{\EdgeLatency}{%
\addplot[azul,thick,mark=*] coordinates {(1000,85.642) (5000,469.275) (10000,539.841) (25000,1479.439) (50000,3115.039)};
\addlegendentry{Q1}
\addplot[verde,thick,mark=square*] coordinates {(1000,160.624) (5000,115.311) (10000,220.581) (25000,329.356) (50000,313.254)};
\addlegendentry{Q2}
}
\newcommand{\EdgeThroughput}{%
\addplot[azul,thick,mark=*] coordinates {(1,1.800) (2,3.600) (5,7.600)};
\addlegendentry{Q1}
\addplot[verde,thick,mark=square*] coordinates {(1,6.400) (2,11.600) (5,21.200)};
\addlegendentry{Q2}
}
\newcommand{\EdgeTemperature}{%
\addplot[verde,thick] coordinates {(0.000,47.774) (0.087,47.774) (0.174,48.312) (0.261,48.850) (0.348,47.774) (0.436,48.312) (0.522,48.312) (0.614,48.312) (0.697,48.312) (0.784,48.312) (0.871,49.388) (0.958,48.312) (1.045,49.388) (1.132,49.388) (1.219,49.388) (1.306,48.312) (1.393,49.388) (1.480,49.388) (1.568,49.388) (1.654,49.388) (1.742,48.850) (1.829,49.388) (1.916,49.388) (2.023,48.312) (2.091,49.388) (2.178,49.388) (2.265,49.388) (2.352,49.388) (2.439,50.464) (2.527,49.388) (2.616,49.388) (2.702,49.926) (2.787,50.464) (2.874,49.388) (2.961,50.464) (3.048,50.464) (3.135,50.464) (3.222,50.464) (3.311,49.926) (3.396,50.464) (3.483,51.002) (3.570,50.464) (3.657,50.464) (3.744,51.002) (3.831,49.926) (3.918,51.002) (4.007,51.540) (4.092,51.002) (4.179,51.002) (4.266,51.540) (4.353,51.540) (4.440,51.002) (4.527,51.540) (4.614,51.540) (4.701,51.540) (4.788,51.540) (4.875,52.078) (4.962,52.078) (5.049,52.616) (5.136,52.616) (5.231,51.540) (5.311,51.540) (5.397,52.078) (5.486,51.540) (5.571,53.154) (5.658,52.616) (5.745,52.616) (5.833,53.692) (5.920,52.616) (6.007,52.616) (6.094,53.154) (6.181,54.230) (6.268,54.768) (6.355,53.154) (6.443,56.382) (6.533,55.306)};
\addplot[verde,thick] coordinates {(8.345,49.388) (8.432,52.616) (8.524,49.388) (8.606,51.002) (8.698,51.540) (8.798,52.616) (8.885,52.616) (8.973,53.154) (9.060,52.616) (9.147,53.154) (9.234,53.692) (9.321,54.768) (9.408,53.692) (9.496,55.844) (9.583,54.768) (9.671,56.920) (9.758,55.306) (9.845,55.844) (9.933,55.306) (10.021,55.844) (10.108,54.768) (10.195,56.382) (10.282,54.768) (10.372,54.768) (10.461,55.306) (10.548,55.306) (10.634,56.920) (10.726,56.920) (10.809,55.844) (10.896,56.382) (10.983,56.920) (11.071,56.920) (11.158,55.306) (11.246,58.534) (11.333,57.996) (11.453,57.458) (11.507,57.996) (11.595,57.458) (11.682,57.996) (11.769,57.996) (11.856,57.996) (11.950,57.996) (12.030,57.458) (12.117,58.534) (12.205,59.072) (12.291,57.996) (12.379,57.996) (12.466,58.534) (12.556,59.072) (12.640,59.072) (12.727,57.458) (12.816,62.300) (12.906,62.300) (12.995,62.300) (13.085,62.838) (13.172,64.452) (13.261,64.452) (13.350,64.452) (13.442,66.604) (13.530,65.528) (13.622,64.452) (13.707,66.604) (13.798,68.218) (13.887,66.604) (13.979,65.528) (14.066,68.218) (14.155,65.528) (14.255,67.142)};
}
\newcommand{\EdgeRam}{%
\addplot[verde,thick] coordinates {(0.000,251.492) (0.087,260.367) (0.174,257.648) (0.261,253.176) (0.348,257.230) (0.436,254.008) (0.522,247.562) (0.614,252.906) (0.697,256.039) (0.784,254.246) (0.871,252.113) (0.958,253.762) (1.045,258.059) (1.132,254.082) (1.219,258.426) (1.306,252.945) (1.393,244.422) (1.480,257.047) (1.568,257.484) (1.654,259.504) (1.742,257.789) (1.829,256.539) (1.916,258.422) (2.023,252.820) (2.091,255.840) (2.178,254.961) (2.265,253.949) (2.352,245.246) (2.439,255.090) (2.527,253.074) (2.616,247.523) (2.702,252.020) (2.787,253.051) (2.874,256.418) (2.961,251.227) (3.048,254.777) (3.135,257.508) (3.222,252.016) (3.311,247.676) (3.396,246.293) (3.483,244.844) (3.570,246.996) (3.657,247.574) (3.744,245.223) (3.831,235.734) (3.918,248.691) (4.007,249.535) (4.092,252.273) (4.179,250.059) (4.266,248.352) (4.353,249.297) (4.440,248.680) (4.527,243.727) (4.614,245.109) (4.701,235.457) (4.788,221.207) (4.875,209.148) (4.962,203.176) (5.049,219.852) (5.136,228.254) (5.231,235.816) (5.311,231.973) (5.397,222.441) (5.486,234.910) (5.571,213.016) (5.658,214.727) (5.745,225.070) (5.833,203.785) (5.920,220.285) (6.007,230.203) (6.094,209.594) (6.181,209.141) (6.268,164.113) (6.355,208.641) (6.443,204.305) (6.533,170.766)};
\addplot[verde,thick] coordinates {(8.345,199.969) (8.432,197.875) (8.524,197.289) (8.606,200.371) (8.698,204.207) (8.798,197.535) (8.885,211.137) (8.973,198.434) (9.060,222.516) (9.147,204.391) (9.234,207.945) (9.321,202.336) (9.408,231.398) (9.496,235.266) (9.583,225.980) (9.671,236.863) (9.758,199.574) (9.845,200.234) (9.933,200.398) (10.021,197.602) (10.108,244.066) (10.195,201.633) (10.282,226.281) (10.372,193.961) (10.461,199.680) (10.548,228.895) (10.634,230.352) (10.726,218.699) (10.809,207.727) (10.896,194.980) (10.983,196.719) (11.071,200.141) (11.158,231.402) (11.246,229.047) (11.333,216.609) (11.453,225.520) (11.507,207.902) (11.595,207.906) (11.682,217.723) (11.769,215.590) (11.856,226.152) (11.950,213.352) (12.030,205.730) (12.117,219.758) (12.205,208.504) (12.291,232.508) (12.379,218.078) (12.466,210.473) (12.556,224.629) (12.640,218.766) (12.727,221.484) (12.816,206.910) (12.906,192.277) (12.995,210.289) (13.085,203.922) (13.172,213.293) (13.261,197.730) (13.350,220.852) (13.442,205.793) (13.530,182.695) (13.622,209.133) (13.707,207.250) (13.798,210.633) (13.887,211.004) (13.979,233.004) (14.066,209.430) (14.155,211.090) (14.255,206.473)};
}
\newcommand{\EdgeFrequency}{%
\addplot[verde,thick] coordinates {(0.000,1000.000) (0.087,1000.000) (0.174,1000.000) (0.261,1000.000) (0.348,700.000) (0.436,1000.000) (0.522,700.000) (0.614,1000.000) (0.697,1000.000) (0.784,1000.000) (0.871,1000.000) (0.958,1000.000) (1.045,1000.000) (1.132,1000.000) (1.219,1000.000) (1.306,1000.000) (1.393,1000.000) (1.480,1000.000) (1.568,1000.000) (1.654,1000.000) (1.742,1000.000) (1.829,1000.000) (1.916,1000.000) (2.023,1000.000) (2.091,1000.000) (2.178,1000.000) (2.265,1000.000) (2.352,1000.000) (2.439,1000.000) (2.527,1000.000) (2.616,1000.000) (2.702,1000.000) (2.787,1000.000) (2.874,1000.000) (2.961,1000.000) (3.048,1000.000) (3.135,1000.000) (3.222,1000.000) (3.311,1000.000) (3.396,1000.000) (3.483,1000.000) (3.570,1000.000) (3.657,1000.000) (3.744,1000.000) (3.831,1000.000) (3.918,1000.000) (4.007,1000.000) (4.092,1000.000) (4.179,1000.000) (4.266,1000.000) (4.353,1000.000) (4.440,1000.000) (4.527,1000.000) (4.614,1000.000) (4.701,1000.000) (4.788,1000.000) (4.875,1000.000) (4.962,1000.000) (5.049,1000.000) (5.136,1000.000) (5.231,1000.000) (5.311,1000.000) (5.397,1000.000) (5.486,1000.000) (5.571,1000.000) (5.658,1000.000) (5.745,1000.000) (5.833,1000.000) (5.920,1000.000) (6.007,1000.000) (6.094,1000.000) (6.181,1000.000) (6.268,1000.000) (6.355,1000.000) (6.443,1000.000) (6.533,1000.000)};
\addplot[verde,thick] coordinates {(8.345,1000.000) (8.432,1000.000) (8.524,1000.000) (8.606,1000.000) (8.698,1000.000) (8.798,1000.000) (8.885,1000.000) (8.973,1000.000) (9.060,1000.000) (9.147,1000.000) (9.234,1000.000) (9.321,1000.000) (9.408,1000.000) (9.496,1000.000) (9.583,1000.000) (9.671,1000.000) (9.758,1000.000) (9.845,1000.000) (9.933,1000.000) (10.021,1000.000) (10.108,1000.000) (10.195,1000.000) (10.282,1000.000) (10.372,1000.000) (10.461,1000.000) (10.548,1000.000) (10.634,1000.000) (10.726,1000.000) (10.809,1000.000) (10.896,1000.000) (10.983,1000.000) (11.071,1000.000) (11.158,1000.000) (11.246,1000.000) (11.333,1000.000) (11.453,1000.000) (11.507,1000.000) (11.595,1000.000) (11.682,1000.000) (11.769,1000.000) (11.856,1000.000) (11.950,1000.000) (12.030,1000.000) (12.117,1000.000) (12.205,1000.000) (12.291,1000.000) (12.379,1000.000) (12.466,1000.000) (12.556,1000.000) (12.640,1000.000) (12.727,1000.000) (12.816,1000.000) (12.906,1000.000) (12.995,1000.000) (13.085,1000.000) (13.172,1000.000) (13.261,1000.000) (13.350,1000.000) (13.442,1000.000) (13.530,1000.000) (13.622,1000.000) (13.707,1000.000) (13.798,1000.000) (13.887,1000.000) (13.979,1000.000) (14.066,1000.000) (14.155,1000.000) (14.255,1000.000)};
}


\begin{document}
\begin{titlepage}
\thispagestyle{fancy}
\vspace*{21mm}
{\color{verde}\large\bfseries COMPUTACIÓN DE BORDE \hfill ESTUDIO DE CASO\par}
\vspace{10mm}
{\fontsize{34}{39}\selectfont\bfseries\color{azul}
¿Hasta dónde llega una\\ Raspberry Pi con JSON?\par}
\vspace{6mm}
{\Large Latencia, memoria y temperatura de PostgreSQL JSONB en una Zero 2 W\par}
\vspace{12mm}
{\large Proyecto EDGE-DB · IT2iD\par}
{\normalsize Bases de Datos en la Nube · Actividad 2026.C3.EVA03\par}
\vspace{13mm}
\noindent\fcolorbox{veladura}{veladura}{\begin{minipage}{0.943\linewidth}
\vspace{4pt}
\textbf{\color{azul}El hallazgo.} En la placa probada, PostgreSQL 17.12 completó dos consultas de agregación JSONB con hasta 50,000 pedidos y un cliente. Con 10,000 pedidos se ensayaron también dos y cinco clientes. La latencia P95 fue de \EdgeDato{q1_50000_1_p95}\ ms para Q1 y \EdgeDato{q2_50000_1_p95}\ ms para Q2 a 50,000 pedidos. Son medianas de cinco repeticiones y corresponden al recorrido completo observado desde la laptop por SSH.
\vspace{4pt}
\end{minipage}}\par
\vspace{10mm}
\entrada{La pregunta detrás de la cifra}
¿Puede un equipo pequeño resumir datos cerca de donde se producen, sin mandar cada evento a un servidor de nube? Acercar el procesamiento al origen reduce la dependencia de una infraestructura remota en algunos escenarios \cite{nist}. El reto es cuánto trabajo cabe en la memoria, el almacenamiento y el margen térmico de una sola placa.

En este experimento, un pedido es un documento con fecha, cliente, pago y lista de artículos. Se eligió JSONB, la representación interna de PostgreSQL para datos JSON que facilita procesar documentos sin volver a interpretar el texto completo en cada operación \cite{pgjson}. El estudio mide hasta dónde llegó esta configuración y cómo se comportó mientras trabajaba.

\vfill
{\small \textbf{Alcance.} Una placa, una microSD y una condición térmica. La evaluación de MongoDB es documental. Los datos de los pedidos son sintéticos. Las mediciones provienen del hardware real.}
\end{titlepage}
\setcounter{page}{2}

\section{Diseño y método}
La Raspberry Pi Zero 2 W Rev 1.0 incorpora un Cortex-A53 de cuatro núcleos y 512 MB nominales \cite{pi}. Se utilizó Raspberry Pi OS Lite de 64 bits, Debian 13.7, kernel \clave{6.18.50+rpt-rpi-v8}, y una microSD ADATA de 64 GB UHS-I, Class 10, V10. El usuario confirmó el montaje sin disipador. La fuente y el cable quedaron sin identificar y no se midieron ambiente ni energía. La propuesta original contempla dos placas \cite{playbook}. Este estudio registra como desviación el uso de una sola.

\entrada{Factibilidad y validez de las consultas}
PostgreSQL se instaló mediante APT. Se cargaron los cinco conjuntos en una base aislada, se compararon ambas salidas con un oráculo Python y se comprobó la persistencia de 1,000 pedidos y su Q1 tras reiniciar el clúster. En MongoDB Community 8.0, los binarios ARM64 oficiales requieren ARMv8.2-A o posterior \cite{mongo}. El manual del Cortex-A53 documenta Armv8-A \cite{arm}. Por esa incompatibilidad documentada, MongoDB quedó fuera de la ejecución y del contraste de rendimiento.

\textbf{Q1} desanida artículos con \clave{jsonb\_array\_elements} \cite{pgfunctions}, conserva precios mayores a 50 y agrupa subtotales por país y categoría. \textbf{Q2} toma pedidos completados entre el 1 de marzo inclusive y el 1 de diciembre exclusivo de 2025. Devuelve los 1,000 mayores por total de pago, con desempate por identificador ascendente. La tabla usa una clave primaria sobre el identificador y una columna JSONB. Se conservó la configuración del motor, con \clave{shared\_buffers} de 128 MiB y \clave{work\_mem} de 4 MiB.

\medskip
\begin{center}
\small
\begin{tabular}{@{}rrrrr@{}}
\toprule
Pedidos & Artículos & Califican en Q2 & JSONL (MiB) & En PostgreSQL (MiB)\\
\midrule
\EdgeDatasetRows
\bottomrule
\end{tabular}
\par\smallskip
\textit{Tabla 1.} Prefijos de un generador con semilla 202603. El tamaño almacenado incluye la clave primaria. Los hashes SHA-256 están en el manifiesto.
\end{center}

\entrada{La medición desde la laptop}
El piloto ejecutó ambas consultas tres veces en 1,000 pedidos con un cliente. El plan inicial fijó \EdgeDato{warm}\ s de calentamiento, \EdgeDato{window}\ s de medición y \EdgeDato{cool}\ s de espera entre intentos. Se usó caché caliente. Se midieron cinco tamaños con un cliente y, en 10,000 pedidos, también dos y cinco clientes. Se varió el orden Q1/Q2 con semilla fija dentro de cada repetición.

En una corrida de Q1 a 50,000 pedidos el resultado válido llegó fuera de la ventana inicial. La enmienda conservó esos intentos y reinició las dos celdas de 50,000 pedidos. Amplió Q1 a 20 s y separó el calentamiento de la medición en la continuación. Q2 y las celdas de concurrencia conservaron 5 s. El cambio de ventana y orden de sesión limita la comparación directa entre celdas.

La laptop envió las consultas por SSH a clientes persistentes de \clave{psql} y recibió todas las filas. El reloj empieza al enviar la consulta y termina al recibir el marcador de fin. Por eso la cifra incluye la red, el procesamiento y la transferencia, además del tiempo del motor. Cinco pings previos dieron un RTT medio de 6 ms y un rango de 4 a 9 ms, una fotografía inicial de la red. La validación de las filas ocurre después del cronómetro. La telemetría se guardó a una cadencia próxima a 1 Hz.

\newpage
\section{Qué muestran las corridas}
\entrada{1. Análisis exploratorio}
Se registraron \EdgeDato{attempts}\ intentos: \EdgeDato{pilot_runs}\ piloto, \EdgeDato{final_runs}\ incorporados a la matriz final, \EdgeDato{superseded_completed}\ completos bajo el plan inicial de 50,000 pedidos y \EdgeDato{failed}\ fallido por la ventana. Las \EdgeDato{total_operations}\ operaciones finales aceptadas coincidieron con el oráculo. Las salidas fuera de ventana quedaron marcadas. El identificador de arranque \EdgeDato{boot_check}\ al inicio y al cierre. La búsqueda de OOM en el kernel quedó \EdgeDato{kernel_oom}.

El ensayo previo a instalar el motor registró 121 muestras en 602.208 s y 267.3 MiB de RAM disponible en promedio. El reposo con PostgreSQL aportó \EdgeDato{idle_samples}\ muestras en \EdgeDato{idle_duration}\ s, con \EdgeDato{idle_ram}\ MiB disponibles en promedio. Son capturas independientes, de distinta duración y estado del sistema. Su diferencia describe las dos sesiones y no estima por sí sola el costo causal del motor.

\entrada{2. Análisis descriptivo}
\textbf{P95} es el tiempo bajo el que terminó aproximadamente el 95\,\% de las operaciones de una corrida. P99 describe una cola más extrema. Los percentiles de cola ayudan a ver demoras que el promedio oculta \cite{tail}. Se usó rango más cercano. En \EdgeDato{low_p99}\ corridas finales hubo menos de 100 operaciones válidas, por lo que P99 puede representar una de las lecturas más altas y resulta inestable.

\begin{center}
\small
\begin{tabular}{@{}rrlrrrr@{}}
\toprule
Pedidos & Clientes & Carga & Ops. mín. & P95 [P25, P75] (ms) & P99 (ms) & Consultas/s\\
\midrule
\EdgeRows
\bottomrule
\end{tabular}
\par\smallskip
\textit{Tabla 2.} Medianas de cinco repeticiones por celda. Los corchetes muestran cuartiles 25 y 75 del P95. «Ops. mín.» es el menor número de operaciones aceptadas por repetición.
\end{center}

\begin{center}
\begin{tikzpicture}
\begin{axis}[width=.94\linewidth,height=58mm,xmode=log,
  xtick={1000,5000,10000,25000,50000},xticklabels={1k,5k,10k,25k,50k},
  ylabel={P95 (ms)},xlabel={Pedidos · un cliente},grid=major,
  legend pos=north west,legend style={font=\footnotesize}]
\EdgeLatency
\end{axis}
\end{tikzpicture}
\par\smallskip
{\small \textit{Figura 1.} Mediana de P95 frente al volumen. El eje horizontal usa escala logarítmica. Los cuartiles están en la tabla 2.}
\end{center}

\newpage
\section{Diagnóstico y límites}
\entrada{3. Qué ayuda a explicar la degradación}
Al pasar de 1,000 a 50,000 pedidos con un cliente, el P95 de Q1 pasó de \EdgeDato{q1_1000_1_p95}\ a \EdgeDato{q1_50000_1_p95}\ ms y el de Q2 de \EdgeDato{q2_1000_1_p95}\ a \EdgeDato{q2_50000_1_p95}\ ms. Cada consulta tiene un costo distinto. Q1 visita los artículos anidados y agrupa los que pasan el filtro. Q2 devuelve hasta 1,000 filas y parte de su tiempo de respuesta puede corresponder al envío de resultados.

A 50,000 pedidos, Q1 registró un ordenamiento externo (\textit{\EdgeDato{50000-q1-sort}}) y \EdgeDato{50000-q1-temp}\ bloques temporales escritos. Q2 usó \textit{\EdgeDato{50000-q2-sort}} y escribió \EdgeDato{50000-q2-temp}\ bloques temporales. El tiempo del motor fue \EdgeDato{50000-q1-ms}\ y \EdgeDato{50000-q2-ms}\ ms respectivamente. Estos planes de \clave{EXPLAIN ANALYZE} se tomaron fuera del benchmark y omiten el envío de las filas del resultado \cite{pgexplain}. El ordenamiento externo da evidencia directa de trabajo temporal en disco en Q1. Una comprobación adicional sin compilación JIT dio \EdgeDato{50000-q1-jit-off-ms}\ ms en una sola invocación, insuficiente para aislar el costo del JIT.

\begin{center}
\begin{tikzpicture}
\begin{axis}[width=.49\linewidth,height=53mm,xtick={1,2,5},xmin=.7,xmax=5.3,
  ylabel={Consultas/s},xlabel={Clientes · 10,000 pedidos},grid=major,
  legend pos=north west,legend style={font=\scriptsize},tick label style={font=\scriptsize},label style={font=\footnotesize}]
\EdgeThroughput
\end{axis}
\end{tikzpicture}
\hfill
\begin{tikzpicture}
\begin{axis}[width=.47\linewidth,height=53mm,
  ylabel={Temperatura del SoC (°C)},xlabel={Minutos desde el piloto},grid=major,
  tick label style={font=\scriptsize},label style={font=\footnotesize}]
\EdgeTemperature
\end{axis}
\end{tikzpicture}
\par\smallskip
{\small \textit{Figuras 2 y 3.} Throughput mediano a 10,000 pedidos (izq.) y traza térmica de las dos etapas (der., punto cada cinco muestras). El espacio entre líneas corresponde a la pausa de la enmienda.}
\end{center}

En 10,000 pedidos, Q1 pasó de \EdgeDato{q1_10000_1_qps}\ a \EdgeDato{q1_10000_5_qps}\ consultas/s al subir de uno a cinco clientes. En Q2 pasó de \EdgeDato{q2_10000_1_qps}\ a \EdgeDato{q2_10000_5_qps}. A la vez, el P95 de Q2 subió de \EdgeDato{q2_10000_1_p95}\ a \EdgeDato{q2_10000_5_p95}\ ms. En esta configuración, el aumento de capacidad vino acompañado de una respuesta de cola más lenta.

En las ventanas finales, la RAM disponible bajó hasta \EdgeDato{run_ram_min}\ MiB y la suma de RSS de procesos PostgreSQL llegó a \EdgeDato{run_rss_max}\ MiB. Esta suma puede contar memoria compartida más de una vez. La CPU global observada llegó a \EdgeDato{run_cpu_max}\,\%, el swap hasta \EdgeDato{run_swap_max}\ MiB y el SoC a \EdgeDato{run_temp_max}\ °C. Las máscaras observadas fueron \clave{\EdgeDato{flags}}. Los máximos independientes de actividad microSD por ventana fueron \EdgeDato{read_mib}\ MiB leídos y \EdgeDato{write_mib}\ MiB escritos. La cadencia de muestreo, la caché y la ruta Wi-Fi impiden atribuir las colas a un único recurso.

\entrada{Qué pudimos observar sobre el voltaje}
En un ensayo separado, el firmware rechazó la consulta del canal de entrada de 5 V. Registramos 64 muestras del riel interno VC4 durante reposo, carga Q1 a 10,000 pedidos y recuperación. Sus medianas fueron 1.2500, 1.2625 y 1.2500 V respectivamente \cite{firmware}. La familia Zero carece del circuito de aviso de bajo voltaje documentado para otros modelos \cite{power}. Por tanto, incluso con máscara \clave{0x0}, el voltaje real de entrada y la energía quedan sin medir.

\newpage
\section{Qué se puede afirmar}
El máximo volumen probado con éxito fue de 50,000 pedidos a un cliente. En 10,000 pedidos se observaron hasta cinco clientes. Es una frontera de cobertura: configura un punto reproducible para esta unidad, pero deja abierta la capacidad por encima de esos niveles o con concurrencia a 50,000 pedidos. Se midió el SoC en una única condición térmica. Quedaron fuera del análisis causal el efecto de un disipador y el consumo eléctrico.

Como recomendación de ingeniería, conviene definir qué latencia de cola tolera la aplicación. El rango intercuartílico aquí describe variación entre corridas, mientras que las diferencias entre placas y sesiones exigirían replicación en esos niveles \cite{benchmark}. Futuras sesiones pueden alargar las ventanas, medir RTT continuamente y ejecutar un diseño cruzado seguro con y sin disipador. Una comparación específica del transporte puede conservar la consulta y separar el tiempo del motor y de la respuesta completa.

\entrada{Cómo seguir la evidencia}
\begin{tabularx}{\linewidth}{@{}lX@{}}
\toprule
Ruta del proyecto & Evidencia\\
\midrule
\clave{hardware/} & Inventario de la placa y características físicas declaradas.\\
\clave{dataset/} & Generador, archivos JSONL y manifiesto con semilla y hashes.\\
\clave{workloads/} & Consultas PostgreSQL verificadas y equivalentes propuestas para MongoDB.\\
\clave{results/raw/} & Línea base, ejecución por operación, telemetría, planes y logs.\\
\clave{results/processed/} & Resumen por celda, cálculos y planes normalizados.\\
\clave{protocol/} & Diseño, bitácora, prompt maestro y auditoría de IA.\\
\bottomrule
\end{tabularx}

\medskip
\textbf{Uso de IA.} Codex asistió con código, ejecución, análisis y redacción. Los logs y las mediciones son salidas de la placa y del cliente, y el oráculo es un cálculo independiente en Python. El repositorio contiene la trazabilidad de esa asistencia y los puntos pendientes de revisión humana.

\begin{thebibliography}{10}
\small\raggedright
\bibitem{nist} M. Iorga y otros. \textit{Fog Computing Conceptual Model}. NIST SP 500-325, 2018. \href{https://doi.org/10.6028/NIST.SP.500-325}{DOI: 10.6028/NIST.SP.500-325}.
\bibitem{pgjson} PostgreSQL Global Development Group. \textit{PostgreSQL 17: JSON Types}. \href{https://www.postgresql.org/docs/17/datatype-json.html}{Documentación oficial}.
\bibitem{pi} Raspberry Pi Ltd. \textit{Zero models}. \href{https://www.raspberrypi.com/documentation/computers/zero.html}{Especificaciones oficiales}.
\bibitem{playbook} Universidad Politécnica de Chiapas. \textit{Caracterización de bases de datos JSON en computación de borde}. Playbook EDGE-DB v1.3, Bases de Datos en la Nube, 2026. Documento de la actividad.
\bibitem{mongo} MongoDB Inc. \textit{MongoDB 8.0: Production Notes}, apartado ARM64. \href{https://www.mongodb.com/docs/v8.0/administration/production-notes/#arm64}{Requisitos oficiales}.
\bibitem{arm} Arm. \textit{Cortex-A53 MPCore Processor Technical Reference Manual}, sección 1.2.1. \href{https://documentation-service.arm.com/static/6040c321ee937942ba301626}{Manual de arquitectura}.
\bibitem{pgfunctions} PostgreSQL Global Development Group. \textit{PostgreSQL 17: JSON Functions and Operators}. \href{https://www.postgresql.org/docs/17/functions-json.html}{Funciones de JSONB}.
\bibitem{tail} J. Dean y L. A. Barroso. \textit{The Tail at Scale}. Communications of the ACM 56(2), 2013. \href{https://research.google/pubs/the-tail-at-scale/}{Publicación de los autores}.
\bibitem{pgexplain} PostgreSQL Global Development Group. \textit{PostgreSQL 17: EXPLAIN}. \href{https://www.postgresql.org/docs/17/sql-explain.html}{Referencia de SQL}.
\bibitem{benchmark} T. Kalibera y R. E. Jones. \textit{Rigorous Benchmarking in Reasonable Time}. ISMM, 2013. \href{https://doi.org/10.1145/2464157.2464160}{DOI: 10.1145/2464157.2464160}.
\bibitem{firmware} Raspberry Pi Ltd. \textit{Raspberry Pi OS: vcgencmd}. \href{https://www.raspberrypi.com/documentation/computers/os.html#measure_volts-block}{Voltajes internos y comandos de firmware}.
\bibitem{power} Raspberry Pi Ltd. \textit{Raspberry Pi computer hardware: Power supply warnings}. \href{https://www.raspberrypi.com/documentation/computers/raspberry-pi.html#power-supply-warnings}{Detección de bajo voltaje por familia}.
\end{thebibliography}

\newpage
\appendix
\section{Recursos durante el ensayo}
La RAM disponible y la frecuencia de CPU 0 forman parte de \EdgeDato{trace_samples}\ observaciones instrumentadas. Las dos series incluyen intervalos entre corridas. Una separación entre líneas indica la pausa entre etapas.

\begin{center}
\begin{tikzpicture}
\begin{axis}[width=.96\linewidth,height=62mm,ymin=64,grid=major,
  xlabel={Minutos desde el piloto},ylabel={RAM disponible (MiB)}]
\EdgeRam
\end{axis}
\end{tikzpicture}
\par\smallskip
{\small \textit{Figura A1.} RAM disponible, muestreada en la Pi y registrada en la laptop. Se muestra un punto cada cinco muestras. El protocolo fijó 64 MiB disponibles como umbral preventivo.}
\end{center}

\begin{center}
\begin{tikzpicture}
\begin{axis}[width=.96\linewidth,height=62mm,ymin=560,ymax=1050,grid=major,
  xlabel={Minutos desde el piloto},ylabel={Frecuencia de CPU 0 (MHz)}]
\EdgeFrequency
\end{axis}
\end{tikzpicture}
\par\smallskip
{\small \textit{Figura A2.} Lectura de \texttt{scaling\_cur\_freq} de CPU 0. Se muestra un punto cada cinco muestras. La frecuencia reportada complementa la máscara de limitación y la temperatura del SoC.}
\end{center}
\end{document}
